%% kapa: exam=1 round=31 year=2020 period=1 form=A subject=civil_law subject_ko=민법 seq=6 number=6
%% source: Q-Net 감정평가사 기출문제 (cst00309) artlSeq=5207219 file=extracted/1008/5207219/A형/A형/2020년도 제31회 감정평가사 1차 시험_1교시_A형.pdf
%% conversion: kordoc 4.24.51 -> tools/gen1.py (auto); review=proofread-sonnet; flags=none
\begin{question}{6}
법인에 관한 설명으로 옳지 않은 것은? (다툼이 있으면 판례에 따름)
\begin{choices}
  \item 법인은 설립등기를 함으로써 성립한다.
  \item 어느 사단법인과 다른 사단법인의 동일 여부는, 다른 사정이 없으면 사원의 동일 여부를 기준으로 결정된다.
  \item 법인의 대표자는 그 명칭이나 직위 여하가 아니라 법인등기를 기준으로 엄격하게 확정하여야 한다.
  \item 행위의 외형상 직무행위로 인정할 수 있으면, 대표자 개인의 이익을 위한 것이거나 법령에 위반한 것이라도 직무에 관한 행위이다.
  \item 대표자의 행위가 직무에 관한 것이 아님을 알았거나 중대한 과실로 모른 피해자는 법인에 손해배상책임을 물을 수 없다.
\end{choices}
\end{question}
\answer{3}
